%%%%%%%%%%%%%%%%%
% This is an sample CV template created using altacv.cls
% (v1.7.4, 30 July 2025) written by LianTze Lim (liantze@gmail.com). Compiles with pdfLaTeX, XeLaTeX and LuaLaTeX.
%
%% It may be distributed and/or modified under the
%% conditions of the LaTeX Project Public License, either version 1.3
%% of this license or (at your option) any later version.
%% The latest version of this license is in
%%    http://www.latex-project.org/lppl.txt
%% and version 1.3 or later is part of all distributions of LaTeX
%% version 2003/12/01 or later.
%%%%%%%%%%%%%%%%

%% Use the "normalphoto" option if you want a normal photo instead of cropped to a circle
% \documentclass[10pt,a4paper,withhyper,normalphoto]{altacv}

\documentclass[10pt,a4paper,withhyper]{../altacv}
%% AltaCV uses the fontawesome5 and simpleicons packages.
%% See http://texdoc.net/pkg/fontawesome5 and http://texdoc.net/pkg/simpleicons for full list of symbols.

% Change the page layout if you need to
\geometry{left=1.25cm,right=1.25cm,top=1.5cm,bottom=1.5cm,columnsep=1.2cm}

% The paracol package lets you typeset columns of text in parallel
\usepackage{paracol}

% Change the font if you want to, depending on whether
% you're using pdflatex or xelatex/lualatex
% WHEN COMPILING WITH XELATEX PLEASE USE
% xelatex -shell-escape -output-driver="xdvipdfmx -z 0" sample.tex
\iftutex
  % If using xelatex or lualatex:
  \setmainfont{Roboto Slab}
  \setsansfont{Lato}
  \renewcommand{\familydefault}{\sfdefault}
\else
  % If using pdflatex:
  \usepackage[rm]{roboto}
  \usepackage[defaultsans]{lato}
  % \usepackage{sourcesanspro}
  \renewcommand{\familydefault}{\sfdefault}
\fi

% Change the colours if you want to
\definecolor{SlateGrey}{HTML}{2E2E2E}
\definecolor{LightGrey}{HTML}{666666}
\definecolor{DarkPastelRed}{HTML}{450808}
\definecolor{PastelRed}{HTML}{8F0D0D}
\definecolor{GoldenEarth}{HTML}{E7D192}
\colorlet{name}{black}
\colorlet{tagline}{PastelRed}
\colorlet{heading}{DarkPastelRed}
\colorlet{headingrule}{GoldenEarth}
\colorlet{subheading}{PastelRed}
\colorlet{accent}{PastelRed}
\colorlet{emphasis}{SlateGrey}
\colorlet{body}{LightGrey}

% Change some fonts, if necessary
\renewcommand{\namefont}{\Huge\rmfamily\bfseries}
\renewcommand{\personalinfofont}{\footnotesize}
\renewcommand{\cvsectionfont}{\LARGE\rmfamily\bfseries}
\renewcommand{\cvsubsectionfont}{\large\bfseries}


% Change the bullets for itemize and rating marker
% for \cvskill if you want to
\renewcommand{\cvItemMarker}{{\small\textbullet}}
\renewcommand{\cvRatingMarker}{\faCircle}
% ...and the markers for the date/location for \cvevent
% \renewcommand{\cvDateMarker}{\faCalendar*[regular]}
% \renewcommand{\cvLocationMarker}{\faMapMarker*}


% If your CV/résumé is in a language other than English,
% then you probably want to change these so that when you
% copy-paste from the PDF or run pdftotext, the location
% and date marker icons for \cvevent will paste as correct
% translations. For example Spanish:
% \renewcommand{\locationname}{Ubicación}
% \renewcommand{\datename}{Fecha}


%% Use (and optionally edit if necessary) this .tex if you
%% want to use an author-year reference style like APA(6)
%% for your publication list
% \input{pubs-authoryear}

%% Use (and optionally edit if necessary) this .tex if you
%% want an originally numerical reference style like IEEE
%% for your publication list
\usepackage[backend=biber,style=ieee,sorting=ydnt,defernumbers=true]{biblatex}
%% For removing numbering entirely when using a numeric style
\setlength{\bibhang}{1.25em}
\DeclareFieldFormat{labelnumberwidth}{\makebox[\bibhang][l]{\itemmarker}}
\setlength{\biblabelsep}{0pt}
\defbibheading{pubtype}{\cvsubsection{#1}}
\renewcommand{\bibsetup}{\vspace*{-\baselineskip}}
\AtEveryBibitem{%
  \iffieldundef{doi}{}{\clearfield{url}}%
}


%% sample.bib contains your publications
\addbibresource{sample.bib}

\begin{document}
\name{Guanlin Wang}
\tagline{Software \& IT Engineer}
%% You can add multiple photos on the left or right
\photoR{2.8cm}{../../images/gl_photo}
% \photoL{2.5cm}{Yacht_High,Suitcase_High}

\personalinfo{%
  % Not all of these are required!
  \email{g.wang@mail.de}
  \phone{+49 151 56312173}
  \mailaddress{Alfred-Messel-Weg 38, 64287, Darmstadt}
  \printinfo{\faLinkedin}{Guanlin Wang}[https://www.linkedin.com/in/guanlin-wang-737354206]
}

\makecvheader
%% Depending on your tastes, you may want to make fonts of itemize environments slightly smaller
% \AtBeginEnvironment{itemize}{\small}

%% Set the left/right column width ratio to 6:4.
\columnratio{0.3}

% Start a 2-column paracol. Both the left and right columns will automatically
% break across pages if things get too long.
\begin{paracol}{2}

  \cvsection{FÄHIGKEITEN}

  % Don't overuse these \cvtag boxes — they're just eye-candies and not essential. If something doesn't fit on a single line, it probably works better as part of an itemized list (probably inlined itemized list), or just as a comma-separated list of strengths.
  \cvsubsection{Systeme \& Deployment}
  \cvtag{Linux}

  \vspace{0.5em}

  \cvtag{Ubuntu / WSL}

  \vspace{0.5em}

  \cvtag{Docker}

  \vspace{0.5em}

  \cvtag{SSH}

  \vspace{0.5em}
  
  \cvtag{Shell / Bash}
  
  \vspace{0.5em}

  \cvtag{Batch}

  \vspace{0.5em}

  \divider\smallskip

  \cvsubsection{Softwareentwicklung}
  
  \cvtag{Python}
  
  \vspace{0.5em}

  \cvtag{C++}
  
  \vspace{0.5em}

  \cvtag{C\#}
  
  \vspace{0.5em}

  \cvtag{.NET}
  
  \vspace{0.5em}

  \cvtag{REST API}
  
  \vspace{0.5em}

  \cvtag{FastAPI}
  
  \vspace{0.5em}

  \cvtag{ASP.NET Core}
  
  \vspace{0.5em}

  \divider\smallskip

  \cvsubsection{Datenbanken \& Integration}
  
  \cvtag{Microsoft SQL Server}
  
  \vspace{0.5em}

  \cvtag{SQL}
  
  \vspace{0.5em}

  \cvtag{System Integration}

  \vspace{0.5em}

  \divider\smallskip

  \cvsubsection{AI \& Bildverarbeitung}
  
  \cvtag{Image Processing}
  
  \vspace{0.5em}  
  
  \cvtag{OCR}
  
  \vspace{0.5em}  
  
  \cvtag{Computer Vision}
  
  \vspace{0.5em}  
  
  \cvtag{OpenCV}
  
  \vspace{0.5em}  
  
  \cvtag{PyTorch}

  \vspace{0.5em}  \vspace{10em}

  \cvsection{SPRACHEN}

  \cvtag{Chinesisch (Muttersprache)}

  \vspace{0.5em}

  \cvtag{Deutsch (C1)}

  \vspace{0.5em}

  \cvtag{Englisch (B2)}


  %% Yeah I didn't spend too much time making all the
  %% spacing consistent... sorry. Use \smallskip, \medskip,
  %% \bigskip, \vspace etc to make adjustments.
  \medskip

  % use ONLY \newpage if you want to force a page break for
  % ONLY the current column
  \newpage

  %% Switch to the right column. This will now automatically move to the second
  %% page if the content is too long.
  \switchcolumn
  \cvsection{BERUFSERFAHRUNG}

  \cvevent{Software | AI \& Computer Vision Entwickler}{Innomatik AG}{06/2021 -- Heute}{Bensheim}

  \vspace{1em}

  % linux system
  \cvsubsubsection{Linux, Entwicklungsumgebungen \& Automatisierung}

  \vspace{1.0em}

  \begin{itemize}
    \item Nutzung von \textbf{Ubuntu unter WSL} als zentrale Entwicklungsumgebung sowie Einrichtung und Konfiguration projektspezifischer Software- und AI-Umgebungen
    \item Nutzung von \textbf{SSH} für die Entwicklung und Arbeit auf internen Linux-Servern
    \item Erstellung von \textbf{Shell- und Batch-Skripten} zur Automatisierung von Programmstarts, Parameterübergabe, Datei- und Pfadprüfung sowie zur Einrichtung benötigter Laufzeitkomponenten
  \end{itemize}

  \vspace{1.0em}

  \textbf{Technologien}: Linux, Ubuntu, WSL, Shell, Bash, Batch, SSH, Git
  \vspace{1.5em}

  % deployment
  \cvsubsubsection{Deployment, Containerisierung \& Troubleshooting}

  \vspace{1.0em}

  \begin{itemize}
    \item Bereitstellung eigener Anwendungen und REST-Services auf lokalen Entwicklungsrechnern, internen Servern sowie innerhalb bestehender Unternehmenssysteme
    \item Containerisierung und Ausführung verschiedener Anwendungen und AI-Services mit \textbf{Docker}
    \item Konfiguration von REST-Endpunkten, Ports und Laufzeitparametern für den internen Zugriff auf bereitgestellte Services
    \item Eigenständige Analyse und Behebung technischer Probleme bei Deployment und Anwendungsausführung, unter anderem bei fehlerhaften URL-Routen, Port-Konflikten, nicht startenden Docker-Containern sowie Anwendungs- und Konfigurationsfehlern
  \end{itemize}

  \vspace{1.0em}

  \textbf{Technologien}: Docker, REST-API, FastAPI, ASP.NET Core, Python, C++, C\#

  \vspace{1.5em}

  % software development
  \cvsubsubsection{Softwareentwicklung \& Systemintegration}

  \vspace{1.0em}

  \begin{itemize}
    \item Entwicklung modularer Softwarekomponenten und Verarbeitungspipelines mit Python, C++ und C\# sowie Bereitstellung über REST-Schnittstellen
    \item Entwicklung von Services mit FastAPI und ASP.NET Core WebAPI sowie Integration in bestehende .NET Blazor-Server-Systeme
    \item Integration von ONNX-basierten AI-Komponenten in C++- und C\#-Anwendungen sowie Entwicklung eigenständiger Desktop- und Konsolenanwendungen für interne technische Workflows
  \end{itemize}

  \vspace{1.0em}

  \textbf{Technologien}: Python, C++, C\#, .NET, ASP.NET Core, Blazor Server, FastAPI, REST-API, ONNX

  \vspace{1.5em}

  \newpage

  % database
  \cvsubsubsection{Datenbanken \& Softwareentwicklung}

  \vspace{1.0em}

  \begin{itemize}
    \item Regelmäßige Arbeit mit \textbf{Microsoft SQL Server} im Rahmen der Entwicklung bestehender Unternehmenssoftware
    \item Analyse von Datenbankstrukturen, Datentypen und Beziehungen sowie deren Abbildung in Softwaremodellen und Lese-/Schreibprozessen
    \item Erstellung und Anpassung von Testdaten zur Validierung neuer und bestehender Softwarefunktionen
  \end{itemize}

  \vspace{1.0em}

  \textbf{Technologien}: Microsoft SQL Server, SQL, .NET

  \vspace{1.5em}

  % computer vision
  \cvsubsubsection{OCR \& industrielle Bildverarbeitung}

  \vspace{1.0em}

  \begin{itemize}
    \item Entwicklung einer modularen \textbf{OCR-Pipeline} für reale Bilder und technische Zeichnungen mit Text Detection und Text Recognition
    \item Aufbau einer vollständigen Trainings- und Verarbeitungspipeline einschließlich Datenvorbereitung, Annotation, Modelltraining sowie Bildvor- und Nachverarbeitung
    \item Bereitstellung der Anwendungen als lokale Software, REST-Service und containerisierte Anwendung
  \end{itemize}

  \vspace{1.0em}

  \textbf{Technologien}: Python, OpenCV, PyTorch, PaddleOCR, ONNX, FastAPI, Docker

  \vspace{1.5em}


  \cvsection{AUSBILDUNG}

  \cvevent{M.Sc. Informatik im Bauwesen}{Technische Universität Darmstadt}{04/2018 -- 03/2021}{Darmstadt}

  \vspace{1.0em}

  \cvsubsection{Schwerpunkte}

  \vspace{1.0em}

  \begin{itemize}
    \item Revit-Plugin-Entwicklung mit \textbf{.NET Framework}
    \item BIM-Datenbanken (CRUD) mit \textbf{SQLite}
    \item BIM-Projektmanagement mit \textbf{MS Project}
  \end{itemize}

  \vspace{1.0em}

  \cvsubsection{Wissenschaftliche Hilfskraft (HiWi)}

  \vspace{1.0em}

  SCOPE (Semantic Construction Project Engineering)
  \begin{itemize}
    \item Entwicklung eines webbasierten ViewCube mit Three.js
    \item Ontologiemodellierung der IFC-Hierarchie mit Protégé
  \end{itemize}

  \vspace{1.0em}

  \divider

  \cvevent{B.Sc. Brückenbau}{Chan'an Universität}{09/2012 -- 06/2016}{Xi'an, China}

  % \divider

  \medskip


\end{paracol}


\end{document}
